\subsubsection*{Basic program: thread count and thread hierarchy}
\inputif{q2_basic.tex}
\inputif{q2_basic_bs.tex}

Tables~\ref{tab:q2basic} and~\ref{tab:q2basicbs} show the basic program for
$m=n=2^{14}$. The best configuration is $t=$\,\qtwobasicbestt{} threads in blocks of
\qtwobasicbestbs{} (\qtwobasicbestms~ms).
\begin{itemize}
\item \textbf{Up to $2^{10}$ threads} the time halves every time $t$ doubles, and it
  flattens out between $2^{11}$ and $2^{12}$.
  There are too few threads to hide the DRAM latency: with $t\le 2^{10}$ there
  are at most 32 warps on 16 SMs, so most SMs have one or two warps, and each
  warp stalls for hundreds of cycles on every load.
\item \textbf{$2^{12}$--$2^{13}$ threads}: the time flattens at about 16.6--17~ms,
  i.e.\ roughly 64~GB/s, which is only about 37\% of peak. A warp load touches 32 different
  cache lines (one per thread, $4mn/t$ bytes apart). We think it reaches even
  this only because each thread walks sequentially through its chunk, so after the first
  miss the next 31 loads of the same thread hit the same 128-byte line in L1.
\item \textbf{Block size matters a lot beyond $2^{13}$ threads}. With blocksize 32 the time stays at 16--17~ms
  up to $2^{17}$ threads, while every larger block size jumps to 26--32~ms.
  Our explanation is the L1 working set (we did not measure cache hit rates
  with a profiler, so this is only an estimate from the cache sizes). Each resident warp keeps 32 lines
  (4~KB) live in L1. An SM can hold 48 warps (1536 threads), which is a 192~KB
  working set, larger than the 128~KB L1/shared array, so lines are evicted
  before their 32 elements are consumed and are fetched again. With 32-thread
  blocks the limit of 16 blocks per SM (compute capability 8.6) caps
  residency at 16 warps per SM (64~KB), which fits in L1. So here lower occupancy
  seems to help, because it avoids cache thrashing. For $t\le 2^{13}$ there are
  at most 16 warps per SM anyway, so there the block size matters only when there are
  fewer blocks than SMs: with $t=2^{13}$ and 1024-thread blocks, or $t=2^{12}$ and
  512-thread blocks, only 8 blocks exist, so only 8 of the 16 SMs work (30.9 and
  31.1~ms instead of about 17~ms).
\item \textbf{Very large $t$} ($\ge 2^{19}$) is slower again. The chunks get short
  (512 elements and less), and the program has to copy $t$ partial
  sums to the host and add them there. Table~\ref{tab:q2trend} shows this cost:
  versions 2 and 3 differ only in the merge, and the gap grows with $t$.
\end{itemize}

\subsubsection*{Effect of the optimisations}
\inputif{q2_versions.tex}
\inputif{q2_trend.tex}

\begin{itemize}
\item \textbf{Coalescing} is the big step: $2.5\times$ over the best basic
  configuration (16.6~ms $\to$ 6.6~ms). Each warp load
  is now one 128-byte transaction that is fully used, so the program reaches
  about 160~GB/s, which is 92\% of peak. Coalescing needs enough threads
  (Table~\ref{tab:q2trend}). With $t\le 2^{12}$ the coalesced version is
  \emph{slower} than the basic one, because each thread then issues one independent
  load per iteration and there are not enough warps to keep enough requests in
  flight to cover the DRAM latency. The basic version is helped there by its L1 hits.
\item \textbf{Tree reduction} on the GPU removes the copy of $t$ partial sums
  and the serial host merge. At the thread count where coalescing is best the gain is
  small (the merge of $2^{16}$ floats is cheap; at $2^{14}$ threads the global-memory tree
  with its $\log_2(\text{blocksize})$ barrier steps is even slightly slower than the host merge), but the gain grows quickly
  with $t$ ($2^{22}$ threads: 13.1~ms $\to$ 6.7~ms). So with the tree reduction
  the best time no longer depends much on $t$.
\item \textbf{Shared memory} for the tree changes little (1.6--6\%). The
  reduction is $\log_2(\text{blocksize})$ steps over at most 1024 values once
  per block, which is negligible compared with streaming 1~GB from DRAM. Its
  global-memory version mostly hits in L1/L2 anyway. Shared memory helps most
  at large $t$, where there are many blocks and therefore many reductions.
\item \textbf{float4 + ILP + warp shuffle} gives the best time and,
  more importantly, reaches it at only $2^{12}$--$2^{13}$ threads. Each load
  instruction moves 16 bytes per thread, and two independent loads per thread are in
  flight, so fewer warps can still keep the memory system busy. The
  shuffle reduction needs no shared memory and only one \texttt{\_\_syncthreads}.
\end{itemize}
All optimised versions converge to 6.3--6.6~ms $\approx$ 160--171~GB/s; the best (v5,
\qtwovfivems~ms) reaches \qtwovfivegbps~GB/s, i.e.\ 97\% of the \devpeakbandwidthgbs~GB/s peak. Since
the kernel is limited by memory bandwidth, we do not expect further optimisations
to gain much; we think the remaining few percent are lost to things like DRAM
refresh and the start and end of the kernel. The numerical difference from the double-precision CPU result is about
$10^{-4}$--$10^{-3}$ for a result of magnitude $\approx 3300$ (relative error
$10^{-7}$--$10^{-6}$). The difference changes between versions because the order of
the float additions changes; it looks like versions that sum many short partial sums
(large $t$) are more accurate than those with long sequential float accumulations
(small $t$).

\subsubsection*{Best CUDA program vs.\ OpenMP}
\inputif{q2_compare.tex}
\inputif{q2_omp.tex}

The best CUDA program (v5, tuned per size) is \qtwospeedupmin--\qtwospeedupmax$\times$
faster than the best OpenMP program (Table~\ref{tab:q2compare}). For every matrix of
64~MB or more the speedup is about $4.3$--$4.7\times$, and it is the same for
square, wide ($2^{12}\times2^{16}$) and tall ($2^{16}\times2^{12}$)
matrices. Both programs seem to be limited by memory bandwidth, so the speedup is roughly
the ratio of the achieved memory bandwidths: about 170~GB/s on the GPU's GDDR6 against 35--40~GB/s
on the CPU's DDR4 (Table~\ref{tab:q2omp}). The CPU seems to reach its memory
bandwidth limit with only 2--4 threads (2 and 4 threads give almost the same time); more threads only seem to add contention and
synchronisation (the 16-thread runs are slower, and for the smallest size the
fork/join cost of 16 threads dominates). For the smallest matrix
($2^{10}\times2^{10}$, 4~MB) there is no gain (both take about 45~$\mu$s): the run time is
tens of microseconds, so we think kernel-launch and copy latency dominate, while the
CPU's 4~MB working set stays in its 16~MB L3 cache between repetitions.
